\section{What Are Quantum Random Oracles?}\label{sec:defs}
  A classical oracle for a function $f$ is a black-box procedure $\Oracle_f$
    such that, given input $x$, it returns $f(x)$ in such a way that no
    information about $f$, other than that which can be inferred from its
    output, is provided. It is helpful to think of the oracle as something that
    exists beyond your reach, ``in the sky'' so to speak, and with which you can
    only interact by sending inputs and receiving outputs.
    A random oracle, then, is simply a classical function oracle where the
    function $f$ was chosen uniformly at random from the space of all functions
    from $\bin^n$ to $\bin^m$.

  However, the operation $x \longrightarrow f(x)$ is generally not reversible,
    and therefore not \emph{unitary}, as is required of operations on quantum
    states. Therefore, we define \emph{quantum-accessible} oracles to be
    oracles that apply the unitary $\unitary_f$, defined as the operation that
    takes the computational basis states $\ket{x, y}$ to $\ket{x, y \xor f(x)}$.\footnote{
      If any of this terminology confuses you, then now would be a good time
      to read through the quantum primer found in \autoref{sec:primer}.
    } In general 
\hhnote{wip}
    registers, respectively), on input a quantum state $\ket{\psi} = \sum_{i
    \in [2^{n+m}]} \alpha_i \ket{x, y}$, we get:
    \[ 
      \sum_{i \in [2^{n+m}]} \alpha_i \ket{x_i, y_i} 
        \apply{O_f} \sum_{i \in [2^{n+m}]} \alpha_i \ket{x_i, y_i \xor f(x_i)}\, .
    \]
    %\[ 
    %  \sum_{i \in [2^{n}]} \sum_{j \in [2^{m}]} \alpha_i \beta_j \ket{i, j} 
    %    \apply{O_f} \sum_{i \in [2^{n}]} \sum_{j \in [2^{m}]} \alpha_i \beta_j \ket{i, j \xor f(i)}\, .
    %\]

  \noindent
  To see that this operation is unitary, it is enough to observe that it is
    \emph{reversible}: Inputing the same state a second time yields the original state: 
    \[ 
      \sum_{i \in [2^{n+m}]} \alpha_i \ket{x_i, y_i \xor f(x_i)}
        \apply{\Oracle_f} \sum_{i \in [2^{n+m}]} \alpha_i \ket{x_i, y_i \xor
        f(x_i) \xor f(x_i)} = \sum_{i \in [2^{n+m}]} \alpha_i \ket{x_i, y_i}\, .
    \]

  \begin{remark} % Footnote, or remark?
    Classical random oracles typically follow the standard for hash functions in
      being defined from the space of \emph{all} finite-length bit strings
      $\bin^*$ to the space of fixed-length bitstrings $\bin^m$. Technically, this means
      that $f$ not a function, but rather a a collection of random functions
      $f_\ell$, one for each input length $\ell$. However, since quantum random
      oracles must operate by unitary means, we will make life easier for
      ourselves by explicitly
      restricting the input domain of the random function to be $\bin^n$, and then just
      take for granted that any result trivially generalizes to
      collections of oracles with varying input length. 
      This follows from the observation that, even when loading up a hash function on 
      real-world quantum hardware, the input \emph{length} to the hash function can
      not be placed in superposition, since this does not correspond to a
      well-defined unitary operation: Every ``branch'' of a superposition needs to
      be defined over the same number of qubits. %% Should be careful about claiming this.
  \end{remark}

  \noindent
  Sometimes it we will want to provide a quantum algorithm with access to a
    classical oracle.
  
  \hhnote{Next.}
